\chapter*{Preface}
\addcontentsline{toc}{chapter}{Preface}

ORIUS is written as a dissertation-style research manuscript rather than as a
flat dump of experiments, proofs, and appendices.  The goal is not to present
every domain as equally mature.  The goal is to make one universal runtime
safety claim readable without hiding where the evidence is deepest, where it is
bounded, and where it is still incomplete.

The battery row remains the reference witness because it is the only surface
with the full theorem-to-code-to-artifact chain.  Autonomous vehicles and
healthcare remain bounded promoted rows: AV through completed nuPlan replay
artifacts, and Healthcare through retrospective monitoring evidence.  That
asymmetry is intentional and is part of the scientific discipline of the
manuscript rather than a defect to be smoothed over by prose.

\paragraph{How to read this dissertation.}
Readers interested primarily in the systems problem should begin with
Chapter~\ref{ch:introduction}, then Chapter~\ref{ch:ch04_theoretical_foundations}
for the runtime objects and theorem discipline, and then
Chapter~\ref{ch:ch07_universal_validation} for the tiered cross-domain
evidence.  Readers focused on the formal surface should pair
Chapter~\ref{ch:ch04_theoretical_foundations} with
Appendix~\ref{app:proofs}, Appendix~\ref{app:verified-theorems-gap-audit}, and
Chapter~\ref{ch:universality-completeness}.  Readers evaluating reproducibility
and promotion discipline should read the conclusion together with
Appendix~\ref{app:integrated-theorem-surface-register},
Appendix~\ref{app:reproducibility-appendix}, and the manifest-backed tables in
\texttt{reports/publication/}.

\paragraph{What this frontmatter adds.}
The sections that follow are brief on purpose.  They do not widen the claim
surface.  They provide the notation, reading conventions, and evidence framing
needed to understand the dissertation without mistaking universal architecture
for equal empirical closure.
